% Generated draft: requires mathematical and source review.
\documentclass[11pt]{article}
\usepackage[margin=25mm]{geometry}
\usepackage{fontspec}
\setmainfont{Latin Modern Roman}
\usepackage{amsmath,amssymb,mathtools,unicode-math}
\setmathfont{Latin Modern Math}
\usepackage{enumitem,needspace,xurl,hyperref}
\setlength{\parindent}{0pt}
\setlength{\parskip}{5pt}
\setlength{\emergencystretch}{3em}
\newcommand{\N}{\mathbb{N}}
\newcommand{\Z}{\mathbb{Z}}
\newcommand{\Q}{\mathbb{Q}}
\newcommand{\R}{\mathbb{R}}
\newcommand{\C}{\mathbb{C}}
\newcommand{\F}{\mathbb{F}}
\newcommand{\A}{\mathbb{A}}
\newcommand{\PP}{\mathbb{P}}
\newcommand{\E}{\mathbb{E}}
\newcommand{\eps}{\varepsilon}
\newcommand{\dd}{\,\mathrm{d}}
\begin{document}
\section*{UnsolvedMath ID: 30004099}
\textbf{Source problem code:} catalogue entry
\subsection{Definitions and mathematical statement}
Let \(k\) be a field and let \(H\) be a finite-dimensional Hopf algebra over \(k\), with coproduct \(\Delta\) and counit \(\varepsilon\). Write \(\tau(x\otimes y)=y\otimes x\). The algebra \(H\) is noncocommutative if \(\Delta\ne\tau\circ\Delta\). Give \(k\) the trivial left \(H\)-module structure \(h\cdot a=\varepsilon(h)a\). Under the usual trivial-coefficient interpretation, its cohomology algebra is
\[
\mathrm H^{\bullet}(H,k)=\bigoplus_{n\geq 0}\operatorname{Ext}^{n}_{H}(k,k),
\]
with the Yoneda product. Finite generation means that there are finitely many homogeneous elements \(u_1,\ldots,u_r\) whose generated unital \(k\)-subalgebra is all of \(\mathrm H^{\bullet}(H,k)\).

The general question, under this interpretation, is: for every finite-dimensional noncocommutative Hopf \(k\)-algebra \(H\), is \(\mathrm H^{\bullet}(H,k)\) finitely generated as a \(k\)-algebra? The supplied record does not specify the permitted base fields, including whether algebraic closure or a characteristic restriction is intended. Thus this formulation is relative to a base field \(k\); its intended range remains unresolved. The record also does not specify whether a stronger finite-generation property involving cohomology with other coefficients is intended.

For the highlighted special case, let \(\operatorname{char}(k)=p\), where \(p>2\) is prime. A Hopf algebra is called connected when its coradical, the sum of its simple subcoalgebras, is \(k\,1\). Among the connected Hopf algebras of dimension \(p^3\), the supplied record asks whether
\[
\mathrm H^{\bullet}(A_5(0),k)
\quad\text{is finitely generated as a \(k\)-algebra.}
\]
Here \(A_5(0)\) denotes the member bearing that name in the source's family \(A_5(\beta)\). No presentation, Hopf structure, parameter domain, or field hypotheses for that family are supplied, so this named special case cannot be specified intrinsically from the record alone.

The original formulation asks about the remaining \(A_5(\beta)\) family in characteristic \(p>2\), including \(\beta=0\). The supplied curation dated 2026-08-22 reports that the nonzero-parameter cases are settled and that finite-generation results cover connected Hopf algebras of dimension \(p^3\) except \(A_5(0)\) for odd \(p\). This explains the narrower special case in the current statement; it is a supplied literature assessment, not an independently verified mathematical conclusion.
\subsection{Short English statement}
Does every finite-dimensional noncocommutative Hopf algebra have a finitely generated cohomology algebra with trivial coefficients? The supplied record particularly asks about the connected algebra \(A_5(0)\) of dimension \(p^3\) in odd characteristic. The intended base-field hypotheses, the precise meaning of cohomological finite generation, and the definition of \(A_5(0)\) require confirmation. The supplied dated curation labels the problem partially solved.
\subsection{Sources}
\begin{itemize}
\item Local UnsolvedMath registry, permanent ID \(30004099\), source problem code \texttt{catalogue entry}; supplied complete source record from the Oberwolfach Reports Open Problems collection. Repository registry locations are \texttt{data/unsolvedmath/problems.json} and \texttt{catalog/problems.json}.
\item Supplied report citation: \emph{Mini-Workshop: Cohomology of Hopf Algebras and Tensor Categories} (2020), Oberwolfach Reports, pp.\ 663--693, DOI \texttt{10.4171/owr/2019/11}. \url{https://doi.org/10.4171/owr/2019/11}. This is also the supplied literature-source link labeled ``Oberwolfach report.'' The record gives proposed year 2019 and no individual proposer; the citation's 2020 date and the 2019 provenance have not been reconciled.
\item Supplied literature source labeled ``Finite-generation for connected Hopf algebras'': \url{https://doi.org/10.1016/j.jalgebra.2019.02.030}. No authors or fuller bibliographic details are supplied.
\item Supplied literature source labeled ``Erdmann on A5 beta'': \url{https://arxiv.org/abs/1904.04533}. The label supplies the Erdmann attribution but no fuller bibliographic details.
\item Literature assessment date: \texttt{2026-08-22}. The supplied dated curation reports finite generation for the nonzero-parameter \(A_5(\beta)\) cases and for connected Hopf algebras of dimension \(p^3\) apart from \(A_5(0)\) in odd characteristic, leaving the general question unresolved. This assessment and the supplied ``partially solved'' label are not independently verified here. The background expressly requests independent verification of time-sensitive status and cited claims.
\end{itemize}
\end{document}
